% sections/A_derivations.tex -- Appendix A: derivations. SKELETON.
% Each derivation: assumptions, statement, proof or derivation, status tag, and the testable
% prediction it feeds into (with a pointer to the figure that tests it). Sources: the theory notes
% notes/theory_escape.md, notes/theory_saddles.md, notes/theory_regimes.md and the agents' notes.
\section{Derivations}\label{app:derivations}

\placeholder{Skeleton. Full derivations behind the statements of the main text; the short proofs of
\cref{sec:setup} are not repeated.}

\subsection{Gradients and the tangent kernel of a TTN}\label{app:gradients}
\placeholder{Environment formula for $\nabla_{\G{v}}\Loss$; the tangent kernel at initialization and
its scale $\isc^{2(\Kc-1)}$.}
% TODO: grad_{G_v} L = contraction of the residual Sigma(W - W*) with all cores except G_v (the
% environment of v); expected kernel at fan-in init; needed for Sec. 5.1 (lazy time rescaling).

\subsection{Escape from the origin}\label{app:escape}
\placeholder{Spherical flow and radial growth for the degree-$\Kc$ homogeneous correlation
$g(\params)=\ip{\W(\params)}{\Wt}$; alignment time; the first-escape formula of
\cref{sec:s2s-escape}.}
% TODO(theory_escape): near the origin, dynamics ~ gradient flow of g. With d tau = |theta|^{K-2} dt
% the direction follows the spherical gradient flow of g and the radius grows like
% exp(K int g). Result to state: t_1 ~ rho_0^{-(K-2)} [tau_align + K^{K/2-1}/((K-2) s_1)];
% aligned limit a_0^{-(K-2)}/((K-2) s). State assumptions (separation of time scales, nondegenerate
% local maxima of g on the sphere). Status: \derived.

\subsection{Expansion around a learned-component saddle}\label{app:saddles}
\placeholder{Second-order expansion of $\Loss$ in the directions that add a new component; the
locality rule for $k$ and the saddle type.}
% TODO(theory_saddles): lowest-degree escape route; reduced ODE for its amplitude; k-rule (H2) and
% strict-iff-k=2 (H3); gauge projection of the Hessian (orbit tangents of eq:gauge).

\subsection{Correction of a wrong learning order}\label{app:piggyback}
\placeholder{The ``piggyback'' direction and its Hessian eigenvalue $q(\tw{\mathrm{learned}}-
\tw{\mathrm{new}})$ (hypothesis H1b).}
% TODO: hand derivation for the N = 4 binary tree with odeco teacher, balanced learned amplitudes
% q = s^{1/K}; generalize.

\subsection{Gradient-descent stability limit}\label{app:gd-stability}
\placeholder{Sharpness at the balanced global minimum and the predicted $\lr_{\max}$.}
% TODO: PLAN.md guess eta_max ~ 2 / (K s_1^{2 - 2/K}): verify (for one balanced mode the Hessian
% of 1/2 (s - prod_v q_v)^2 at q_v = s^{1/K} is s^{2-2/K} 1 1^T, top eigenvalue K s^{2-2/K}) or
% correct it for the full TTN.

\subsection{Regime boundaries}\label{app:boundaries}
\placeholder{Lazy criterion and saddle-to-saddle separation criterion as functions of $\isc$,
$\chi$, $\Kc$ and imbalance; large-bond-dimension scalings.}
% TODO(theory_regimes): Chizat-type criterion; S2S boundary from escape times with separation
% ratio 3; muP-type per-core scalings from the spectral condition.
